\documentclass[11pt]{article}
\usepackage[T1]{fontenc}
\usepackage{lmodern}
\usepackage[margin=1in]{geometry}
\usepackage{amsmath,amssymb,amsthm,booktabs}
\usepackage[hidelinks]{hyperref}
\newtheorem{lemma}{Lemma}
\newtheorem{proposition}[lemma]{Proposition}
\title{Sphere packing and improved scales for the Heilbronn construction}
\author{Ivan Blinov\\\small Research draft prepared with OpenAI Codex assistance}
\date{October 8, 2026}
\begin{document}
\maketitle
\begin{abstract}
An elementary finite form of Behrend's classical sphere construction
replaces the digit packing in the pinned OpenAI Heilbronn manuscript.
Retuning the dimension to \(d=13\) and reducing the auxiliary scales
gives the conditional paper exponent
\[
 k=403^{17},\qquad
 \eta_{\rm tuned}=\frac{1}{1074\cdot403^{17}+7}
       \approx4.77320332233169\cdot10^{-48}.
\]
This is approximately \(3.15045471808690\cdot10^{90}\) times the
previous sphere exponent. The earlier \(d=41\) construction is retained
below as a baseline. Lean checks both complete finite packings, their
binomial counts and capacities, and the new integer scale inequalities.
The full geometric consequence remains conditional on the upstream estimates.
\end{abstract}

\section{Input and scope}
The input is OpenAI math family 191, \emph{A power improvement in the
Heilbronn triangle lower bound}, September 25, 2026, pinned at
\[
 \texttt{fd4aeeb2ee4fc729c18d98444fed42fd0529eeeb}.
\]
All saved upstream files remain unchanged. For the preserved baseline, retain
\[
 d=41,\qquad M=\binom{163}{41},\qquad T=\binom{M}{3}.
\]
The exact \(T\) is recorded in the certificate, independently
recomputed by the verifier, and checked in Lean using the
falling-factorial quotient formula for binomial coefficients.
The original packing has \(k=T^2+1\).
The first local replacement, preserved separately, has \(k=3^{384}\).
The baseline transfer and interface review below use \(d=41\).
Section~\ref{sec:retuned} explicitly changes \(d,T,k\) and the other scales.
This is an application of the classical sphere idea of Behrend
\cite{behrend}, not a new progression-free-set construction.
Stronger variants are known \cite{ehps}; neither optimality nor
novelty of the present application is asserted.

\section{The finite sphere construction}
\begin{lemma}\label{lem:sphere}
Let \(A,m,T\) be positive integers and put
\[
 s=2A,\quad b=4A-1,\quad k=b^m,\quad Q=mA(A-1)/2+1.
\]
If \(TQ\le s^m\), there are \(T\) distinct integers \(u_a\), \(0\le a<T\),
such that
\[
 0\le2u_a\le k-1,\qquad
 u_a+u_{a'}=2u_{a''}\quad\Longleftrightarrow\quad a=a'=a''.
\]
The positions \(i_a=u_a,\ j_a=u_a,\ \ell_a=k-1-2u_a\) therefore lie
in \(\{0,\ldots,k-1\}\), are injective within each row, and satisfy
\[
 i_a+j_{a'}+\ell_{a''}=k-1
 \quad\Longleftrightarrow\quad a=a'=a''.
\]
\end{lemma}
\begin{proof}
For \(x\in\{0,\ldots,s-1\}^m\), set
\[
 X_v=2x_v-(s-1),\qquad E(x)=\sum_{v=0}^{m-1}X_v^2.
\]
Each centered coordinate is odd, has absolute value between \(1\)
and \(2A-1\), and its square is \(1\) modulo \(8\). Hence
\[
 m\le E(x)\le m(2A-1)^2,\qquad E(x)\equiv m\pmod8.
\]
There are at most
\[
 \frac{m((2A-1)^2-1)}8+1=mA(A-1)/2+1=Q
\]
energies. The cube contains exactly \(s^m\) words, so the pigeonhole
principle gives a class with at least \(T\) words. This proves
existence of the needed input rather than assuming it.

Choose an energy class with at least \(T\) words, and take its first
\(T\) words in increasing order of the base-\(s\) index
\(\sum_v x_vs^v\), denoted \(x(a)\). This is the selection used in
the formal existence proof; no least-energy assertion is needed. Define
\[
 u_a=\sum_{v=0}^{m-1}x_v(a)b^v.
\]
Uniqueness of expansion proves injectivity, and
\[
 2u_a\le2(s-1)\sum_{v=0}^{m-1}b^v=b^m-1
\]
because \(2(s-1)=b-1\).

In an equality \(u_a+u_{a'}=2u_{a''}\), each digit on either side
is at most \(b-1\); thus there are no carries. Consequently
\(x(a)+x(a')=2x(a'')\) coordinatewise. Centering gives \(X+Y=2Z\).
These three vectors have equal squared norms, so
\[
 \|X-Y\|^2+4\|Z\|^2=2\|X\|^2+2\|Y\|^2
\]
implies \(X=Y=Z\). The distinct selected words therefore have
equal indices. The converse and the position claims follow directly.
\end{proof}
The selection is finite but enormous; we do not enumerate this sphere
or its \(T\) selected words. The existence theorem does not require
that computation. For \(A=1\), this becomes the binary-digit ternary set.

\section{Baseline parameters and comparison}
Take
\[
 m=29,\quad A=10003,\quad s=20006,\quad b=40011,\quad Q=1450725088.
\]
The exact integer comparison is
\[
 T\cdot1450725088\le20006^{29}.
\]
The right-to-left ratio is approximately \(1.00268267437678\);
only the integer inequality is used in the proof.
The preceding alphabet \(A=10002\) fails this sufficient criterion.

The generator compares dimensions \(3\le m\le284\), selecting the
smallest sufficient \(A\) for each and then comparing \((4A-1)^m\).
The independent checker instead finds the largest \(A\) whose \(k\)
would beat the selected value, and checks that its capacity fails.
Monotonicity of capacity in \(A\) follows since the derivative of
\((2A)^m/(mA(A-1)/2+1)\) has the sign of
\[
 (m-2)A^2-(m-1)A+2
 =(m-3)A(A-1)+(A-1)^2+1>0
\]
for \(m\ge3\), \(A\ge1\). This verifies only the declared bounded
comparison under a sufficient criterion, not unrestricted optimality
or exact maximization of sphere populations.

\begin{center}
\begin{tabular}{lll}
\toprule
Packing & \(k\) & Conditional paper exponent, approximately\\
\midrule
Original & \(T^2+1\) & \(3.17571403643291\cdot10^{-236}\)\\
Ternary & \(3^{384}\) & \(2.68582462243153\cdot10^{-188}\)\\
Sphere & \(40011^{29}\) & \(1.51508393214748\cdot10^{-138}\)\\
Retuned, \(d=13\) & \(403^{17}\) & \(4.77320332233169\cdot10^{-48}\)\\
\bottomrule
\end{tabular}
\end{center}
The sphere-to-ternary exponent ratio is between \(5\cdot10^{49}\)
and \(6\cdot10^{49}\). Lean checks this, the capacity inequality,
and the equality of the supplied \(T\) literal with its binomial formula.

\section{Baseline conditional geometric transfer}
\begin{proposition}
Assume the norm, lattice, orbit, auxiliary-set, and lifting/counting
estimates of the pinned manuscript in their stated generality.
There are constants \(c>0,n_0\) such that every \(n\ge n_0\) admits
\(n\) points in the unit square with all triangle areas at least
\[
 c n^{-2+\eta_{\rm sph}},\qquad
 \eta_{\rm sph}=2/(45435\cdot40011^{29}+16).
\]
\end{proposition}
\begin{proof}
Apply Lemma~\ref{lem:sphere} to the unchanged \(T\) norm-determinant
summands. Retain
\[
 L=r^{10},\quad100k^2L^3<B\le200k^2L^3,\quad
 h=B^k,\quad\tau=\lfloor B^{k-1}/2\rfloor
\]
with \(B\) prime. Conditional on a column's label, draw each of its
\(3k\) digits independently among the \(L/r\) positive integers with
the prescribed residue. Equal position numbers in different rows
do not identify their values or their draws.

In the determinant polynomial
\[
 F(Y)=\det\left(\sum_{v=0}^{k-1}c_{iv}^{(j)}Y^v\right)_{i,j=1}^3
     =\sum_t\gamma_tY^t,
\]
the coefficient \(\gamma_{k-1}\) reduces modulo \(r\) to precisely
the retained norm polynomial: undesignated positions vanish, and
the matching identity leaves the original \(T\) summands.
It is nonzero for distinct labels. For every coefficient,
\[
 |\gamma_t|\le6k^2L^3.
\]
The contribution below degree \(k-1\) is at most
\(3(B^{k-1}-1)/50\). The truncation through that degree represents
the determinant modulo \(h\), has absolute value below \(h/2\),
and, for distinct labels, exceeds \(\tau\). Thus the determinant
obstruction is preserved.

The conditional smaller-Smith-divisor moment permits arbitrary
prescribed residues and retains
\[
 \theta=B(r/L)^4\le200k^2r^{-6}=o(1),\qquad
 {\mathbb E}[D\mid\text{labels}]\le2
\]
for large \(r\). Keep \(h^{100}<q\le2h^{100}\), \(N=(hq)^{10}\), and
\(n_r=\lfloor r\sqrt{N^3/\tau}\rfloor\). The original estimates give
\[
 \frac{{\mathbb E}Z}{n_r}\ll
 n_r(hq)^6N^{-3}
 +n_r^2(\log(2N))^2\frac{h}{N^3r^{41}}=o(1).
\]
The terms are bounded by \(r\tau^{-1/2}(hq)^{-9}=o(1)\) and
\(O_k((\log r)^2r^{-9})\), respectively. Deletion leaves
\[
 n_r^2\Delta(P_r)\ge r^2/64,\quad n_r\asymp_k r^\alpha,\quad
 \alpha=1+30\left(15\cdot101k-\frac{k-1}{2}\right)=45435k+16.
\]
This gives exponent \(2/\alpha\) on the constructed sizes.
The retained prime-interval and subset argument extends it to all
sufficiently large \(n\), using two-sided power bounds rather than
assuming monotonicity of \(n_r\).
\end{proof}

\section{Review of the baseline upstream interfaces}
These hand checks concern compatibility and specific cancellations.
They do not constitute independent certification of every upstream
lemma, and none of these lattice or geometric estimates is formalized here.

\paragraph{Norms and conditional digits.}
Odd \(d\) makes the norm polynomial alternating. Grouping distinct
monomial indices into determinant summands uses no division by \(3!\),
so characteristic three is allowed. The divisor-moment proof fixes
the first row and column. Its unit pivot prescribes four remaining
entries modulo \(B^j\); their independent probabilities are at most
\((r/L)^j\), giving the fourth power. Both arguments retain their
hypotheses with the sphere positions.

\paragraph{Nonzero determinant counts (source Section 2).}
Plane-lattice counts use radius at least the second successive minimum;
otherwise centered short vectors lie on a line. In the anisotropic
short-line case, the dyadic normal count and three column counts give
\[
 U^4Z^2\frac{R_1^2}{UZ}\frac{R_2^2}{UZ}\frac{R_3}{U}
 =R_1^2R_2^2R_3U\le(R_1R_2R_3)^2.
\]
Passing to the sublattice generated by short independent vectors
has bounded index; scaling coordinates yields an injective map into
integer matrices. A nonzero determinant itself supplies three
independent rows of bounded length, as required for that reduction.

\paragraph{Zero determinants (source Section 7).}
Distinct projected columns force rational rank two and a primitive
null vector \(x\) whose coordinates are all nonzero. The plane lattice
has covolume \(J=I|x|/g(x)\), with \(g(x)^3/I\le h^2\).
The affine-independent residue case excludes \(|x|\le H=h^2\);
its dyadic radii are therefore in the range needed for the \(g^3\)
moment. For duplicate residue columns let \(\delta=e_i-e_j\),
and let \(C_q\) denote the matrix of auxiliary residues.
The three remaining losses relative to \(N^6/I^2\), per shell, are
\[
\begin{array}{ll}
 \bar x\notin{\mathbb F}_q\delta:
   & h^2q^{-3}(q^3/s_{\rm aux})^2\ll h^{26}/q,\\
 \bar x\in{\mathbb F}_q\delta,\ \operatorname{rank}C_q\ge2:
   & h^2q^{-2}(q^3/s_{\rm aux})\ll h^{14}/q,\\
 \bar x\in{\mathbb F}_q\delta,\ \operatorname{rank}C_q\le1:
   & h^2q^{-4}(q^3/s_{\rm aux})^2\ll h^{26}/q^2 .
\end{array}
\]
Here \(s_{\rm aux}\) is the auxiliary cap size, not the digit alphabet.
In the first case column equality imposes an index-\(q\) condition
on each row. In the exceptional normal class a nonzero coordinate
of \(x\) is divisible by \(q\), ensuring the radius needed for its
residue-class count. In rank at most one, summing the \(q+1\) lines
of the residue plane yields the additional \(O(q^{-2})\) factor
for row triples. All losses are bounded for the retained \(q>h^{100}\).

\paragraph{Probability normalization.}
The lifting factor, orbit bound, and weighted determinant count combine as
\[
 \frac{h^9}{N^9}\frac{D^3E^2}{h^8}
       \frac{N^6(\log N)^2}{D^2E^2}
 =\frac{hD(\log N)^2}{N^3}.
\]
The conditional \(D\)-moment then averages on the label-collision
event. The sphere radius does not enter any of these estimates.
The reviewed interfaces show no incompatibility with the replacement.

\section{Retuning the dimension and scales}\label{sec:retuned}
The digit moment only requires a constant bound \(\theta<1\).
The proof of the auxiliary zero-determinant estimate also permits much
smaller \(H,q\) than the original choices. We now use
\[
\begin{gathered}
 d=13,\qquad M=\binom{51}{13}=476260169700,\\
 T=\binom{M}{3}=18004519550937520342031694999844900,\\
 m=17,\quad A=101,\quad s=202,\quad b=403,\quad Q=85851,\\
 TQ\le202^{17},\qquad k=403^{17}.
\end{gathered}
\]
Lemma~\ref{lem:sphere} supplies the complete packing.
The original norm decomposition applies because \(d\) is still odd.
It uses \(M=\binom{4d-1}{d}\) monomials and \(T=\binom M3\)
determinant summands; no new norm identity is assumed.
The same bounded sphere comparison over \(3\le m\le284\) gives
the displayed \(m,A\) for this smaller \(T\).

\begin{proposition}\label{prop:tuned}
Assume the pinned norm, lattice, orbit, auxiliary-set, and counting
arguments in the general forms used below. There are constants \(c>0,n_0\)
such that every \(n\ge n_0\) admits \(n\) points in the unit square
whose triangles all have area at least
\[
 c n^{-2+\eta_{\rm tuned}},\qquad
 \eta_{\rm tuned}=\frac{1}{1074\cdot403^{17}+7}.
\]
\end{proposition}
\begin{proof}
For a sufficiently large odd prime \(r\), take
\[
\begin{gathered}
 L=400k^2r^4,\qquad100k^2L^3<B\le200k^2L^3,\\
 h=B^k,\quad\tau=\lfloor B^{k-1}/2\rfloor,\quad H=2h,\\
 h^{14}<q\le2h^{14},\qquad N=(hq)^4 ,
\end{gathered}
\]
where \(B,q\) are prime. Thus \(B\asymp_k r^{12}\);
both prime intervals are covered by the retained Bertrand lemma.

\paragraph{Digits and the conditional moment.}
Every prescribed residue still has exactly \(L/r=400k^2r^3\)
positive digit choices, and \(L<B\). The coefficient and carry proof
uses only the matching identity and the unchanged bound
\(\lvert\gamma_u\rvert\le6k^2L^3\).
For every fixed triple of labels, including repeated labels,
\[
 \theta=B(r/L)^4\le\frac{200k^2r^4}{L}=\frac12,\qquad
 {\mathbb E}[D\mid\text{labels}]\le2.
\]
This uniform constant is sufficient; decay of \(\theta\) is not needed.
Label collisions have probability at most \(3r^{-13}\).

\paragraph{Auxiliary set and zero determinants.}
For \(h\ge3\), \(h^{14}\ge8000h^2=2000H^2\), so the cap and
restricted-shift construction retain their hypotheses.
The cap size satisfies
\[
 s_{\rm aux}\ge\frac{q^2}{2000^3H^6},\qquad
 q^3/s_{\rm aux}\ll qH^6=64qh^6.
\]
Affinely independent residues exclude \(|x|\le H\), so every contributing
dyadic shell has radius \(R>H/2=h\). This is exactly the range
\(R\ge h\) of the unchanged \(g(x)^3\) moment estimate.
For the other three cases, the losses relative to \(N^6/I^2\) are
\[
\begin{array}{ll}
 \bar x\notin{\mathbb F}_q\delta:
   & h^2q^{-3}(q^3/s_{\rm aux})^2\ll h^{14}/q,\\
 \bar x\in{\mathbb F}_q\delta,\ \operatorname{rank}C_q\ge2:
   & h^2q^{-2}(q^3/s_{\rm aux})\ll h^8/q,\\
 \bar x\in{\mathbb F}_q\delta,\ \operatorname{rank}C_q\le1:
   & h^2q^{-4}(q^3/s_{\rm aux})^2\ll h^{14}/q^2.
\end{array}
\]
All are bounded because \(q>h^{14}\). They need not tend to zero.
The factors \(64,4096\) and the cap constant are absorbed into
absolute constants.
Also \(q>\tau\), \(q\nmid h\), and \(N\) remains an exact multiple
of \(hq\). The lifting formula is therefore unchanged.
The nonzero determinant count still has
\(\lambda_3(\Lambda)\le E\le h<4N\).
These observations preserve both probability bounds used in deletion.

\paragraph{Sample size and deletion.}
Define
\[
 \rho=1074k+6,\quad
 \gamma=\frac{\rho}{2\rho+1},\quad
 a_r=r^\gamma\sqrt{N^3/\tau},\quad n_r=\lfloor a_r\rfloor .
\]
Here \(0<\gamma<1/2\), with
\(2\gamma-1=-1/(2\rho+1)<0\).
Since \(\tau\ge B^{k-1}/3\) and \(N^{3/2}=(hq)^6\),
\[
 n_r(hq)^6N^{-3}
 \le r^\gamma\tau^{-1/2}
 =O_k(r^{\gamma-6(k-1)})=o(1).
\]
The triple term satisfies
\[
\begin{aligned}
 n_r^2(\log(2N))^2\frac{h}{N^3r^{13}}
 &\le(\log(2N))^2\frac{h}{\tau}r^{2\gamma-13}\\
 &=O_k((\log r)^2r^{2\gamma-1})=o(1).
\end{aligned}
\]
Here \(h/\tau\le3B=O_k(r^{12})\) and \(\log(2N)=O_k(\log r)\).
The expected number of violations is \(o(n_r)\), so the same deletion
argument applied to \(2n_r\) samples leaves \(n_r\) distinct projected
points with minimum triangle area at least \(\tau/(16N^3)\).
Furthermore, \(\tau<h\) and \(N\ge h^4\) give
\(a_r>r^\gamma h^{11/2}\to\infty\). Eventually \(n_r\ge a_r/2\), hence
\[
 n_r^2\Delta(P_r)\ge r^{2\gamma}/64.
\]

\paragraph{The exponent and all cardinalities.}
Set
\[
 \beta=6(1+14)k-\frac{k-1}{2}=\frac{179k+1}{2}.
\]
The bounds for \(q,\tau\) yield
\[
 \sqrt2\,r^\gamma B^\beta<a_r
 \le2^6\sqrt3\,r^\gamma B^\beta.
\]
Consequently \(A_kr^\alpha\le n_r\le D_kr^\alpha\) for fixed positive
constants, where
\[
 \alpha=\gamma+12\beta=\rho+\gamma,\qquad
 \eta_{\rm tuned}=\frac{2\gamma}{\alpha}
 =\frac1{\rho+1}=\frac1{1074k+7}.
\]
Using \(r^{2\gamma}\ge D_k^{-\eta_{\rm tuned}}n_r^{\eta_{\rm tuned}}\)
gives the claimed power on the constructed cardinalities.
For an arbitrary large \(n\), take a prime between
\(\lceil(n/A_k)^{1/\alpha}\rceil\) and twice this integer.
Then \(n\le n_r\le C_kn\) for a fixed \(C_k\). Taking an \(n\)-point
subset proves the result for every sufficiently large \(n\).
\end{proof}

The exponent is approximately \(4.77320332233169\cdot10^{-48}\).
Its ratio to the \(d=41\) sphere exponent lies between
\(3\cdot10^{90}\) and \(4\cdot10^{90}\); the exact ratio is in the
certificate. This compares exponents, not finite-size triangle areas.
In particular the collision decay exponent
\(-1/(2\rho+1)\) is very small: constants and thresholds are not
controlled for practical point counts. No optimality or priority is claimed.

\section{Verification and remaining boundary}
For the current result, run \texttt{python scripts/build\_tuned.py},
then \texttt{python scripts/verify\_tuned.py}. The ternary,
sphere-only, and retuned patches are three alternatives applied
directly to the unchanged manuscript, not successive patches.
The retuned patch updates Sections 3--8, including the full deletion
and interpolation argument.
The earlier sphere checks remain available through
\texttt{python scripts/verify\_sphere.py}.

The verifier independently recomputes binomials by an exact-division
recurrence, the triple count by a cubic expansion, powers by repeated
multiplication, and layer counts by triangular-number increments.
It checks complete small spheres, corrupted certificates, deliberately
invalid layouts, determinant coefficients and carries, native Lean
compilation, patch applicability, and all saved upstream hashes.
The report records exact coverage and evidence hashes.

\texttt{lean/Sphere.lean} now proves the complete finite construction.
The proof chain is:
\begin{enumerate}
\item \texttt{decode\_value} and \texttt{decode\_injective} produce
\(s^m\) distinct digit words from the indices \(0,\ldots,s^m-1\).
\item \texttt{energy\_bound} and \texttt{norm\_energy} prove a
nonnegative layer index \(J(x)\le mA(A-1)/2\) with
\(E(x)=m+8J(x)\). Internally this uses triangular numbers of the
absolute centered digit radii.
\item \texttt{large\_fiber} proves the finite pigeonhole principle
by bounding every fiber and summing its contribution.
\texttt{select\_large\_fiber} selects the first \(T\) indices from
a qualifying fiber and proves their distinctness using list
nonduplication.
\item \texttt{sphere\_family} combines the decoder, energy bound,
and selection. \texttt{packing\_exists} combines this family with
sphere rigidity and carry-free encoding.
\item \texttt{paperT\_binomial} checks
\(\binom{\binom{163}{41}}3\) using the falling-factorial quotient
definition: \(n(n-1)\cdots(n-k+1)/k!\). Recursion is in \(k\), not
in the large upper argument. This uses a local definition, not
Mathlib's binomial API.
\item \texttt{paper\_packing\_binomial} establishes the final
packing with \(T=\binom{\binom{163}{41}}3\), \(k=40011^{29}\),
all position bounds, both row injections, and the matching equivalence.
\end{enumerate}
Thus the formal result no longer assumes the existence of a sufficiently
large equal-energy family. The selected words are defined within the
existence proof; the enormous list is never evaluated during checking.
The proofs use only the standard logical axioms
\texttt{propext}, \texttt{Classical.choice}, and \texttt{Quot.sound}.
The numerical comparison theorems use no axioms.
Finite examples remain separate implementation checks.

\texttt{lean/Tuned.lean} imports this general sphere proof and establishes
\texttt{tuned\_packing} with
\(T=\binom{\binom{51}{13}}3\), \(k=403^{17}\).
It also proves digit divisibility, the inequality \(2Br^4\le L^4\),
the shell threshold \(H/2=h\), the cap-modulus inequality for \(h\ge3\),
and the three integer bounds underlying the zero-determinant cases.
Cleared-denominator inequalities certify the two negative deletion
powers and the identity for \(\eta_{\rm tuned}\); an exact numerical
theorem bounds the gain. These integer facts do not constitute
formal proofs of the probability estimates or limits.

The retuned verifier rebuilds the imported sphere proof before checking
the new file, audits their logical axioms, independently derives the
rational exponent from the scales, and verifies all 16 upstream files.
It rejects four corrupted certificates and five changes that break the
stated sufficient estimates: \(d=11\), \(q\asymp h^{13}\),
\(H=h\), \(N=(hq)^3\), and changing the digit multiplier from \(400\)
to \(399\) while still asserting \(\theta\le1/2\).
These controls show failures of the chosen bounds, not impossibility
of those parameters under a different argument.

The complete geometric theorem, practical point generation, novelty,
priority, and unrestricted optimality remain outside the established
scope. The unchanged upstream Lean comparator uses exponent
\(1/(100000k)\) and a subsequence target; it has not been rerun.

\begin{thebibliography}{9}
\bibitem{upstream} OpenAI, \emph{A power improvement in the Heilbronn
triangle lower bound}, September 25, 2026.
Exact paths and hashes: \texttt{upstream/manifest.json}.
\url{https://github.com/openai/math}.
\bibitem{behrend} F. A. Behrend, \emph{On Sets of Integers Which Contain
No Three Terms in Arithmetical Progression}, PNAS 32 (1946), 331--332.
\url{https://doi.org/10.1073/pnas.32.12.331}.
\bibitem{ehps} C. Elsholtz, Z. Hunter, L. Proske, L. Sauermann,
\emph{Improving Behrend's construction: Sets without arithmetic
progressions in integers and over finite fields}, 2024.
\url{https://arxiv.org/abs/2406.12290}.
\end{thebibliography}
\end{document}
